\chapter{Metodología}

\section{Diagnóstico de la Situación Actual}
Este diagnóstico describe cómo se aprovechan hoy los microdatos de la ENSANUT 2018 y qué características de los datos condicionan el desarrollo. El proceso actual se reconstruyó a partir de la documentación oficial de la encuesta y del trabajo previo de exploración del autor; no se midieron tiempos de ejecución ni se entrevistó a usuarios, por lo que esos aspectos se declaran como limitación.

\subsection{Proceso actual (AS-IS)}
La Tabla~\ref{tab:asis} resume los elementos del proceso y la Figura~\ref{fig:asis} su flujo.

\begin{table}[H]
  \caption{Proceso actual de aprovechamiento de los microdatos de la ENSANUT 2018.}
  \label{tab:asis}
  \tablabase\small
  \begin{tabular}{|p{3.0cm}|p{10.6cm}|}
    \hline
    \enc{Elemento} & \enc{Descripción}\\ \hline
    Usuarios & Investigadores, estudiantes y analistas que consultan la ENSANUT.\\ \hline
    Fuentes & Microdatos públicos del INEC (módulo de salud de la niñez), diccionario de datos y metadatos de la encuesta.\\ \hline
    Herramientas & Tablas y paneles publicados por el INEC; para análisis propios, hojas de cálculo o software estadístico. No se identificó una herramienta específica en la UTI.\\ \hline
    Salidas & Prevalencias por grupo (área, etnia, quintil de ingresos). No hay clasificación individual ni métricas por clase.\\ \hline
    Responsables & INEC (publicación); el analista (procesamiento e interpretación).\\ \hline
    Limitaciones & Códigos numéricos que requieren el diccionario; \num{2017} registros sin variable objetivo; desbalance de clases; procesamiento manual difícil de reproducir.\\ \hline
  \end{tabular}
\end{table}

\begin{figure}[H]
  \centering
  \begin{tikzpicture}[font=\footnotesize, node distance=5mm,
    caja/.style={draw=black!60, rounded corners=2pt, align=center, text width=3.0cm, minimum height=1.3cm, fill=white},
    flecha/.style={-{Stealth[length=2mm]}, black!60}]
    \node[caja] (a) {El INEC publica microdatos y diccionario};
    \node[caja, right=of a] (b) {Descarga y codificación manual};
    \node[caja, right=of b] (c) {Análisis descriptivo (prevalencias)};
    \node[caja, right=of c] (d) {Informe o panel de indicadores};
    \draw[flecha] (a) -- (b); \draw[flecha] (b) -- (c); \draw[flecha] (c) -- (d);
  \end{tikzpicture}
  \caption{Flujo actual (AS-IS) de aprovechamiento de los microdatos de la ENSANUT. Elaboración propia.}
  \label{fig:asis}
\end{figure}

\subsection{Hallazgos sobre los datos}
La exploración inicial de la base produjo los resultados de la Tabla~\ref{tab:datos}. El hallazgo con mayor efecto metodológico fue que el \SI{9.8}{\percent} de los registros no tenía medida la variable objetivo.

\begin{table}[H]
  \caption{Características de la base de datos.}
  \label{tab:datos}
  \tablabase\small
  \begin{tabular}{|p{8.2cm}|r|}
    \hline
    \enc{Característica} & \enc{Valor}\\ \hline
    Registros originales (módulo de salud de la niñez) & \num{20510}\\ \hline
    Variables originales & 350\\ \hline
    Registros sin variable objetivo medida & \num{2017} (\SI{9.8}{\percent})\\ \hline
    Registros del modelo (con variable objetivo) & \num{18493}\\ \hline
    Predictores del modelo & 114\\ \hline
    Casos con desnutrición crónica (clase 1) & \num{4560} (\SI{24.7}{\percent})\\ \hline
    Casos sin desnutrición crónica (clase 0) & \num{13933} (\SI{75.3}{\percent})\\ \hline
    Rango de edad & 0 a 59 meses\\ \hline
    Registros de Tungurahua con variable objetivo & 531 (164 casos; \SI{30.9}{\percent})\\ \hline
    Predictores con valores ausentes en la base original & 84 de 114 (32 con más del \SI{10}{\percent})\\ \hline
  \end{tabular}
\end{table}

\subsection{Síntesis del diagnóstico}
Del diagnóstico se derivan tres necesidades: (i) un procedimiento de preparación reproducible que excluya, y no imponga, la variable objetivo faltante; (ii) una evaluación con métricas por clase que considere el desbalance; y (iii) un medio de consulta (API y tablero) que permita usar el modelo y revisar su desempeño. Estas necesidades sustentan los requerimientos de la sección correspondiente.

\section{Metodología de Desarrollo o Implementación}
\subsection{Enfoque de investigación}
La investigación es aplicada, cuantitativa, no experimental y transversal, porque usa datos secundarios de una sola medición para construir y evaluar un clasificador. El nivel de análisis es de clasificación con componente correlacional; no se buscan relaciones causales.

\subsection{Metodología OSEMN}
OSEMN se ejecutó como proceso de cinco fases (Figura~\ref{fig:osemn}). La Tabla~\ref{tab:osemn} indica, para cada fase, sus entradas, actividades, decisiones, salidas y evidencia. Los scripts y notebooks citados están en el repositorio del proyecto (Anexo~E).

\begin{figure}[H]
  \centering
  \begin{tikzpicture}[font=\footnotesize, node distance=4mm,
    caja/.style={draw=black!60, rounded corners=2pt, align=center, text width=2.3cm, minimum height=1.6cm, fill=white},
    flecha/.style={-{Stealth[length=2mm]}, black!60}]
    \node[caja] (o) {\textbf{Obtener}\\ ENSANUT 2018 (CSV)};
    \node[caja, right=of o] (s) {\textbf{Limpiar}\\ Base de \num{18493} registros};
    \node[caja, right=of s] (e) {\textbf{Explorar}\\ Perfil y control de fuga};
    \node[caja, right=of e] (m) {\textbf{Modelar}\\ Comparación y selección};
    \node[caja, right=of m] (n) {\textbf{Interpretar}\\ Métricas, API y tablero};
    \foreach \a/\b in {o/s,s/e,e/m,m/n} \draw[flecha] (\a) -- (\b);
  \end{tikzpicture}
  \caption{Fases de la metodología OSEMN aplicadas en el proyecto. Elaboración propia.}
  \label{fig:osemn}
\end{figure}

\begingroup
\scriptsize\setlength{\tabcolsep}{3pt}\arrayrulecolor{black!55}\renewcommand{\arraystretch}{1.15}
\begin{longtable}{|p{1.3cm}|p{2.0cm}|p{3.4cm}|p{2.6cm}|p{2.2cm}|p{1.9cm}|}
  \caption{Fases de OSEMN ejecutadas: entradas, actividades, decisiones, salidas y evidencia.}
  \label{tab:osemn}\\ \hline
  \enc{Fase} & \enc{Entradas} & \enc{Actividades} & \enc{Decisiones} & \enc{Salidas} & \enc{Evidencia}\\ \hline
  \endfirsthead
  \tablacont{6}\hline
  \enc{Fase} & \enc{Entradas} & \enc{Actividades} & \enc{Decisiones} & \enc{Salidas} & \enc{Evidencia}\\ \hline
  \endhead
  \hline\endfoot
  Obtener & Microdatos ENSANUT 2018; diccionario y metadatos & Carga del CSV; revisión de estructura y de tipos & Usar solo ENSANUT en el modelo; ENDI solo exploratoria & Base bruta: \num{20510} $\times$ 350 & Notebook de limpieza\\ \hline
  Limpiar & Base bruta & Renombrar variables; convertir valores especiales en nulos; eliminar registros sin variable objetivo; definir 114 predictores & Excluir (no imputar) la variable objetivo faltante & Base del modelo: \num{18493} $\times$ 115 & Script de limpieza\\ \hline
  Explorar & Base del modelo & Distribuciones y prevalencias; comparación Tungurahua y nacional; revisión de fuga por variable & Mantener variables sin señal de fuga (AUC univariado $<$ 0.95) & Perfil de datos; lista de limitaciones & Script de revisión de fuga; notebook EDA\\ \hline
  Modelar & Base del modelo & Partición 80/20 estratificada (semilla 42); preprocesamiento en \textit{Pipeline}; 4 algoritmos con manejo de desbalance; búsqueda de hiperparámetros con validación cruzada; ajuste del umbral & Selección por PR-AUC; umbral que maximiza F1 de la clase 1 (validación cruzada en entrenamiento) & \texttt{modelo\_v2.joblib}; \texttt{metricas\_v2.json} & Script de entrenamiento\\ \hline
  Interpretar & Modelo y particiones & Métricas por clase con intervalos; importancia de variables; API y tablero; pruebas funcionales & Reportar sobreajuste y límites & Resultados (cap. III) y prototipo & Anexos A, C y D\\ \hline
\end{longtable}
\endgroup

\subsection{Protocolo experimental}
La base del modelo se dividió en \num{14794} registros de entrenamiento y \num{3699} de prueba (partición 80/20 estratificada por la variable objetivo, semilla 42). El preprocesamiento se encapsuló en un \textit{Pipeline} de scikit-learn \cite{pedregosa2011scikit}: variables numéricas con imputación por mediana y estandarización, y variables categóricas con imputación por moda y codificación \textit{one-hot}. La estandarización y la codificación se ajustan solo con los datos de entrenamiento. La imputación del \textit{Pipeline}, en cambio, no tiene efecto porque la base del modelo ya no contiene nulos: los faltantes se imputaron antes de la partición (Tabla~\ref{tab:reglas}).

Para el desbalance se usó ponderación de clases (\texttt{class\_weight="balanced"} en los modelos que la admiten y \texttt{scale\_pos\_weight} en XGBoost). Los hiperparámetros del modelo seleccionado se buscaron con búsqueda aleatoria \cite{bergstra2012random} y validación cruzada estratificada de 5 particiones sobre el entrenamiento, con PR-AUC como criterio. El umbral de decisión se eligió como el valor que maximiza el F1 de la clase positiva con predicciones de validación cruzada sobre el entrenamiento. La comparación de los cuatro modelos base se calculó sobre la partición de prueba, lo que puede hacer levemente optimista el PR-AUC reportado; esta limitación se retoma en el Capítulo~III.

\subsection{Reglas de limpieza y preparación}
La Tabla~\ref{tab:reglas} detalla las reglas aplicadas al construir la base del modelo. El código correspondiente está en el Anexo~I.

\begin{table}[H]
  \caption{Reglas de limpieza y preparación de la base.}
  \label{tab:reglas}
  \tablabase\footnotesize
  \begin{tabular}{|c|p{5.0cm}|p{6.9cm}|}
    \hline
    \enc{N.º} & \enc{Regla} & \enc{Justificación y verificación}\\ \hline
    1 & Excluir los registros con variable objetivo sin medir & Se evita fabricar etiquetas. Se excluyeron \num{2017} registros y quedaron \num{18493}\\ \hline
    2 & Alinear la base cruda y la base renombrada por posición & Se verificó que las etiquetas de los \num{18493} registros medidos coinciden al \SI{100}{\percent}; si difiere, el script se detiene\\ \hline
    3 & Convertir valores especiales en nulos & Unifica códigos de no respuesta antes de imputar\\ \hline
    4 & Imputación de valores ausentes & La base del modelo no tiene nulos: se imputaron con mediana o moda en la limpieza inicial, sobre todos los registros y antes de la partición (limitación; Capítulo~III). El imputador del \textit{Pipeline} queda sin efecto\\ \hline
    5 & Separar entrenamiento y prueba antes del modelado & Partición 80/20 estratificada con semilla 42\\ \hline
    6 & Revisar fuga por variable & AUC univariado $\leq$ 0.95 en todas las variables evaluadas\\ \hline
    7 & Exigir que toda columna sea numérica o categórica & Una prueba detiene el proceso si alguna columna quedara fuera del modelo\\ \hline
  \end{tabular}
\end{table}

\subsection{Búsqueda de hiperparámetros}
Solo el algoritmo ganador recibió búsqueda de hiperparámetros. Para Random Forest se muestrearon 40 combinaciones (\textit{RandomizedSearchCV}) con validación cruzada estratificada de 5 particiones y PR-AUC como criterio (Tabla~\ref{tab:grilla}). Se acotó la profundidad y el tamaño de las hojas para reducir el tamaño del modelo (de 170 MB a 28 MB) y el sobreajuste.

\begin{table}[H]
  \caption{Espacio de búsqueda de Random Forest y valores seleccionados.}
  \label{tab:grilla}
  \tablabase\small
  \begin{tabular}{|p{4.6cm}|p{5.6cm}|c|}
    \hline
    \enc{Hiperparámetro} & \enc{Valores explorados} & \enc{Seleccionado}\\ \hline
    \texttt{n\_estimators} & 100, 150, 200, 300 & 300\\ \hline
    \texttt{max\_depth} & 5, 8, 10, 15, 20 & 20\\ \hline
    \texttt{min\_samples\_leaf} & 1, 2, 5, 10 & 10\\ \hline
    \texttt{min\_samples\_split} & 2, 5, 10 & 2\\ \hline
    \texttt{max\_features} & \texttt{sqrt}, \texttt{log2} & \texttt{sqrt}\\ \hline
  \end{tabular}
\end{table}

\subsection{Riesgos del proyecto y mitigación}
\begin{table}[H]
  \caption{Riesgos identificados y medidas aplicadas.}
  \label{tab:riesgos}
  \tablabase\footnotesize
  \begin{tabular}{|p{3.3cm}|p{4.9cm}|p{5.4cm}|}
    \hline
    \enc{Riesgo} & \enc{Efecto posible} & \enc{Medida en el proyecto}\\ \hline
    Fuga de información & Métricas infladas & Exclusión de la variable objetivo faltante, \textit{Pipeline} y revisión univariada; persiste la imputación previa a la partición\\ \hline
    Desbalance de clases & Modelo que ignora la clase 1 & Ponderación de clases, umbral ajustado y métricas por clase\\ \hline
    Sobreajuste & Desempeño real menor al de entrenamiento & Regularización y reporte de la brecha entrenamiento--prueba\\ \hline
    Representatividad territorial & Estimaciones imprecisas para Tungurahua & Intervalos de confianza y advertencia de uso\\ \hline
    Uso indebido como diagnóstico & Decisiones clínicas sin sustento & Aviso de herramienta académica y alcance declarado\\ \hline
    Acceso no autorizado a la API & Uso o abuso del servicio & Pendiente: no hay autenticación (Capítulo~III)\\ \hline
  \end{tabular}
\end{table}



\subsection{Métricas de evaluación}
Con VP, VN, FP y FN los verdaderos positivos, verdaderos negativos, falsos positivos y falsos negativos de la clase 1 (con desnutrición crónica), se calcularon:
\begin{equation}
  \text{Sensibilidad} = \frac{VP}{VP + FN}
  \eqindice{Sensibilidad (recall)}
\end{equation}
\begin{equation}
  \text{Especificidad} = \frac{VN}{VN + FP}
  \eqindice{Especificidad}
\end{equation}
\begin{equation}
  \text{Precisión} = \frac{VP}{VP + FP}
  \eqindice{Precisión}
\end{equation}
\begin{equation}
  F_1 = 2\,\frac{\text{Precisión}\times\text{Sensibilidad}}{\text{Precisión}+\text{Sensibilidad}}
  \eqindice{F1 de la clase positiva}
\end{equation}
\begin{equation}
  \text{Exactitud balanceada} = \frac{\text{Sensibilidad} + \text{Especificidad}}{2}
  \eqindice{Exactitud balanceada}
\end{equation}
Además se reportan el ROC-AUC y el PR-AUC. La clase prioritaria es la 1 (niños con desnutrición crónica) porque omitirla (falso negativo) es el error más costoso en una aplicación de tamizaje; por eso se prioriza la sensibilidad, aceptando una precisión menor.

\section{Técnicas e Instrumentos de Recolección de Información}
La Tabla~\ref{tab:tecnicas} lista las técnicas realmente aplicadas. No se realizaron entrevistas, encuestas ni validación con expertos, por lo que no se presentan resultados de ese tipo.

\begin{table}[H]
  \caption{Técnicas e instrumentos aplicados.}
  \label{tab:tecnicas}
  \tablabase\small
  \begin{tabular}{|p{3.2cm}|p{5.7cm}|p{4.1cm}|}
    \hline
    \enc{Técnica} & \enc{Instrumento} & \enc{Uso en el proyecto}\\ \hline
    Revisión documental & Artículos indexados (IEEE Xplore, ScienceDirect, MDPI, PubMed) según el protocolo del Capítulo~I & Estado del arte y elección de algoritmos\\ \hline
    Análisis de microdatos & Base ENSANUT 2018 (CSV), diccionario y metadatos; \textit{pandas} & Caracterización y limpieza\\ \hline
    Experimentación computacional & Notebooks y scripts en Python; scikit-learn y XGBoost & Entrenamiento y comparación de modelos\\ \hline
    Pruebas funcionales & Casos de prueba sobre la API (\textit{TestClient} de FastAPI) y pruebas automáticas con \textit{pytest} & Verificación del servicio (Anexo~C)\\ \hline
  \end{tabular}
\end{table}

\section{Requerimientos del Proyecto}
Los requerimientos se formularon de modo que cada uno tenga un criterio de aceptación comprobable. Su cumplimiento se verifica en el Capítulo~III.

\subsection{Requerimientos Funcionales}
\begin{table}[H]
  \caption{Requerimientos funcionales y criterios de aceptación.}
  \label{tab:rf}
  \tablabase\footnotesize
  \begin{tabular}{|c|p{4.1cm}|p{6.6cm}|c|}
    \hline
    \enc{ID} & \enc{Requerimiento} & \enc{Criterio de aceptación} & \enc{Prioridad}\\ \hline
    RF1 & Preprocesamiento reproducible & Imputación y escalado ajustados solo con el entrenamiento; resultados idénticos con semilla 42. & Alta\\ \hline
    RF2 & Comparación de modelos & Cuatro algoritmos evaluados con la misma partición y reporte de sensibilidad, especificidad, precisión, F1, exactitud balanceada, ROC-AUC y PR-AUC. & Alta\\ \hline
    RF3 & Predicción individual & \texttt{POST /predict} con las 114 variables devuelve clase, probabilidad y nivel de riesgo; responde HTTP 400 si las variables no coinciden. & Alta\\ \hline
    RF4 & Predicción masiva & \texttt{POST /predict-file} recibe un archivo y devuelve un \texttt{.xlsx} con una fila por registro. & Media\\ \hline
    RF5 & Consulta de métricas & \texttt{GET /dashboard/metricas} devuelve las métricas del archivo \texttt{metricas\_v2.json}. & Alta\\ \hline
    RF6 & Indicadores del tablero & \texttt{GET /dashboard} admite filtro por provincia; \texttt{GET /dashboard/importancia} lista la importancia de variables. & Media\\ \hline
    RF7 & Documentación interactiva & \texttt{GET /docs} muestra la documentación OpenAPI. & Media\\ \hline
    RF8 & Selección de versión del modelo & La variable de entorno \texttt{MODEL\_VERSION} (v1 o v2) define el modelo activo; por defecto v2. & Baja\\ \hline
  \end{tabular}
\end{table}

\subsection{Requerimientos No Funcionales}
\begin{table}[H]
  \caption{Requerimientos no funcionales y criterios de aceptación.}
  \label{tab:rnf}
  \tablabase\footnotesize
  \begin{tabular}{|c|p{2.6cm}|p{7.5cm}|c|}
    \hline
    \enc{ID} & \enc{Atributo} & \enc{Criterio de aceptación} & \enc{Prioridad}\\ \hline
    RNF1 & Validación de entradas & El \SI{100}{\percent} de las solicitudes con variables faltantes o sobrantes se rechaza con HTTP 400. & Alta\\ \hline
    RNF2 & Rendimiento & Latencia mediana de \texttt{POST /predict} $\leq$ 200 ms, medida sobre 100 solicitudes. & Media\\ \hline
    RNF3 & Disponibilidad & \texttt{GET /health} responde HTTP 200 mientras el servicio esté desplegado. & Media\\ \hline
    RNF4 & Seguridad y acceso & Autenticación y control de acceso a la API (no implementada por decisión de alcance; Capítulo~III). & Media\\ \hline
    RNF5 & Reproducibilidad & Semilla fija (42) y versiones de librerías fijadas en \texttt{requirements.txt}. & Alta\\ \hline
    RNF6 & Mantenibilidad & Código separado por módulos y al menos 5 pruebas automáticas aprobadas. & Media\\ \hline
    RNF7 & Uso responsable & La interfaz y los documentos indican que es una herramienta académica exploratoria y no un sistema de diagnóstico. & Alta\\ \hline
  \end{tabular}
\end{table}

\section{Diseño de la Solución}
\subsection{Arquitectura y flujo de datos}
La solución tiene cuatro componentes (Figura~\ref{fig:arquitectura}): (1) el pipeline de datos y modelado, que produce (2) los artefactos del modelo (\texttt{modelo\_v2.joblib} y \texttt{metricas\_v2.json}); (3) la API REST, que carga los artefactos y expone los servicios; y (4) el tablero web, que consume la API. La Figura~\ref{fig:flujo-pred} muestra el flujo de una predicción individual.

\begin{figure}[H]
  \centering
  \begin{tikzpicture}[font=\footnotesize, node distance=6mm,
    caja/.style={draw=black!60, rounded corners=2pt, align=center, text width=3.5cm, minimum height=1.5cm, fill=white},
    flecha/.style={-{Stealth[length=2mm]}, black!60}]
    \node[caja] (d) {\textbf{Datos}\\ ENSANUT 2018 (CSV)};
    \node[caja, right=of d] (p) {\textbf{1. Pipeline de datos y modelado}\\ Python, scikit-learn};
    \node[caja, right=of p] (a) {\textbf{2. Artefactos}\\ \texttt{modelo\_v2.joblib}\\ \texttt{metricas\_v2.json}};
    \node[caja, below=8mm of a] (api) {\textbf{3. API REST}\\ FastAPI};
    \node[caja, left=of api] (w) {\textbf{4. Tablero web}\\ React};
    \node[caja, left=of w] (u) {\textbf{Usuario}\\ navegador o cliente HTTP};
    \draw[flecha] (d) -- (p); \draw[flecha] (p) -- (a); \draw[flecha] (a) -- (api);
    \draw[flecha] (api) -- (w); \draw[flecha] (w) -- (u);
  \end{tikzpicture}
  \caption{Arquitectura y flujo de datos de la solución. Elaboración propia.}
  \label{fig:arquitectura}
\end{figure}

\begin{figure}[H]
  \centering
  \begin{tikzpicture}[font=\footnotesize, node distance=4mm,
    caja/.style={draw=black!60, rounded corners=2pt, align=center, text width=2.2cm, minimum height=1.5cm, fill=white},
    flecha/.style={-{Stealth[length=2mm]}, black!60}]
    \node[caja] (c) {Cliente envía JSON con 114 variables};
    \node[caja, right=of c] (v) {Validación: nombres y cantidad};
    \node[caja, right=of v] (m) {Pipeline: preprocesador y Random Forest};
    \node[caja, right=of m] (u) {Probabilidad y umbral 0.42};
    \node[caja, right=of u] (r) {Respuesta: clase, probabilidad y riesgo};
    \foreach \a/\b in {c/v,v/m,m/u,u/r} \draw[flecha] (\a) -- (\b);
    \node[draw=black!60, rounded corners=2pt, align=center, text width=2.5cm, below=6mm of v, font=\scriptsize] (e) {HTTP 400 si no coinciden};
    \draw[flecha] (v) -- (e);
  \end{tikzpicture}
  \caption{Flujo de una predicción individual (\texttt{POST /predict}). Elaboración propia.}
  \label{fig:flujo-pred}
\end{figure}

\subsection{Modelo de datos}
El modelo usa 114 predictores de la ENSANUT agrupados en nueve bloques (Tabla~\ref{tab:vars}); el listado completo está en el Anexo~F. La variable objetivo es \texttt{desnutricion\_cronica}. No se incluyen la talla ni el peso actuales del niño; sí se incluyen dos variables de diseño muestral (\texttt{factor\_expansion} y \texttt{estrato}), que se discuten como limitación en el Capítulo~III.

\begin{table}[H]
  \caption{Grupos de variables predictoras.}
  \label{tab:vars}
  \tablabase\footnotesize
  \begin{tabular}{|p{4.6cm}|c|p{6.4cm}|}
    \hline
    \enc{Grupo} & \enc{N.º} & \enc{Ejemplos}\\ \hline
    Territorio y demografía & 10 & provincia, área, sexo, edad en meses, etnia y educación de la madre, estrato\\ \hline
    Orden del hijo & 1 & orden del hijo\\ \hline
    Embarazo y control prenatal & 25 & controles prenatales, micronutrientes, exámenes, consejería\\ \hline
    Parto y recién nacido & 19 & tipo de parto, prematuridad, peso y talla al nacer\\ \hline
    Control posparto & 3 & control posparto y lugar\\ \hline
    Controles y consejería del niño & 15 & controles por edad, consejería en lactancia y alimentación\\ \hline
    Lactancia y convivencia & 5 & duración de la lactancia, convivencia con la madre\\ \hline
    Morbilidad y suplementación & 6 & diarrea, infección respiratoria, hierro, vitamina A\\ \hline
    Vacunación & 30 & dosis por carné y lugar de vacunación\\ \hline
    \textbf{Total} & \textbf{114} & \\ \hline
  \end{tabular}
\end{table}

\subsection{Contrato de la API}
La Tabla~\ref{tab:api} presenta los servicios expuestos. Las respuestas de \texttt{/predict} siguen el esquema \texttt{PrediccionResponse}: \texttt{prediccion} (0/1), \texttt{descripcion}, \texttt{probabilidad} y \texttt{riesgo} (BAJO, MEDIO o ALTO).

\begin{table}[H]
  \caption{Servicios de la API REST.}
  \label{tab:api}
  \tablabase\footnotesize
  \begin{tabular}{|p{3.4cm}|c|p{4.1cm}|p{3.8cm}|}
    \hline
    \enc{Ruta} & \enc{Método} & \enc{Entrada} & \enc{Salida y errores}\\ \hline
    \texttt{/health} & GET & Ninguna & \texttt{\{"status":"OK"\}}\\ \hline
    \texttt{/info} & GET & Ninguna & Nombre de la API, versión y modelo activo\\ \hline
    \texttt{/features} & GET & Ninguna & Número y lista de variables del modelo\\ \hline
    \texttt{/predict} & POST & JSON con las 114 variables & Clase, probabilidad y riesgo; HTTP 400 si no coinciden\\ \hline
    \texttt{/predict-file} & POST & Archivo (formulario multipart) & Archivo \texttt{.xlsx} con predicciones\\ \hline
    \texttt{/dashboard} & GET & Provincia opcional (nombre) & Registros de prueba, casos observados y clasificados, sensibilidad y precisión\\ \hline
    \texttt{/dashboard/}\allowbreak\texttt{metricas} & GET & Ninguna & Métricas del modelo (v2)\\ \hline
    \texttt{/dashboard/}\allowbreak\texttt{importancia} & GET & Ninguna & Importancia por variable\\ \hline
    \texttt{/docs} & GET & Ninguna & Documentación interactiva (OpenAPI)\\ \hline
  \end{tabular}
\end{table}

\subsection{Proceso propuesto (TO-BE)}
Frente al proceso actual (Figura~\ref{fig:asis}), el prototipo automatiza la preparación y agrega la clasificación individual y la consulta de métricas (Figura~\ref{fig:tobe}).

\begin{figure}[H]
  \centering
  \begin{tikzpicture}[font=\footnotesize, node distance=5mm,
    caja/.style={draw=black!60, rounded corners=2pt, align=center, text width=2.4cm, minimum height=1.4cm, fill=white},
    flecha/.style={-{Stealth[length=2mm]}, black!60}]
    \node[caja] (a) {Microdatos ENSANUT 2018};
    \node[caja, right=of a] (b) {Pipeline reproducible de limpieza};
    \node[caja, right=of b] (c) {Entrenamiento y evaluación por clase};
    \node[caja, right=of c] (d) {API REST y tablero web};
    \node[caja, below=7mm of d] (e) {Consulta de predicción, métricas e indicadores};
    \draw[flecha] (a) -- (b); \draw[flecha] (b) -- (c); \draw[flecha] (c) -- (d); \draw[flecha] (d) -- (e);
  \end{tikzpicture}
  \caption{Flujo propuesto (TO-BE) con el prototipo. Elaboración propia.}
  \label{fig:tobe}
\end{figure}

\subsection{Casos de uso}
La Figura~\ref{fig:cu} presenta los casos de uso y la Tabla~\ref{tab:cu} los relaciona con los requerimientos.

\begin{figure}[H]
  \centering
  \begin{tikzpicture}[font=\footnotesize, node distance=6mm,
    actor/.style={draw=black!70, rounded corners=1pt, align=center, minimum width=2.6cm, minimum height=1.0cm, fill=black!6},
    cu/.style={draw=black!60, ellipse, align=center, text width=3.4cm, inner sep=1pt, fill=white},
    linea/.style={black!60}]
    \node[actor] (u) {Usuario\\ (analista)};
    \node[actor, below=28mm of u] (d) {Desarrollador};
    \node[cu, right=32mm of u, yshift=12mm] (c1) {CU1. Consultar indicadores por provincia};
    \node[cu, below=4mm of c1] (c2) {CU2. Clasificar un caso individual};
    \node[cu, below=4mm of c2] (c3) {CU3. Clasificar un archivo de casos};
    \node[cu, below=4mm of c3] (c4) {CU4. Consultar métricas del modelo};
    \node[cu, below=4mm of c4] (c5) {CU5. Consultar la documentación de la API};
    \foreach \c in {c1,c2,c3,c4} \draw[linea] (u) -- (\c);
    \draw[linea] (d) -- (c5); \draw[linea] (d) -- (c4);
  \end{tikzpicture}
  \caption{Casos de uso del prototipo. Elaboración propia.}
  \label{fig:cu}
\end{figure}

\begin{table}[H]
  \caption{Casos de uso y requerimientos asociados.}
  \label{tab:cu}
  \tablabase\small
  \begin{tabular}{|c|p{5.6cm}|p{4.2cm}|c|}
    \hline
    \enc{ID} & \enc{Caso de uso} & \enc{Servicio} & \enc{Req.}\\ \hline
    CU1 & Consultar indicadores por provincia & \texttt{GET /dashboard} & RF6\\ \hline
    CU2 & Clasificar un caso individual & \texttt{POST /predict} & RF3\\ \hline
    CU3 & Clasificar un archivo de casos & \texttt{POST /predict-file} & RF4\\ \hline
    CU4 & Consultar métricas del modelo & \texttt{GET /dashboard/metricas} & RF5\\ \hline
    CU5 & Consultar la documentación & \texttt{GET /docs} & RF7\\ \hline
  \end{tabular}
\end{table}



\section{Justificación Técnica}
La Tabla~\ref{tab:justif} compara cada tecnología elegida con una alternativa usando criterios verificables. La selección del algoritmo se justifica con los resultados del Capítulo~III, no con ventajas generales.

\begin{table}[H]
  \caption{Justificación técnica de las decisiones de diseño.}
  \label{tab:justif}
  \tablabase\footnotesize
  \begin{tabular}{|p{2.4cm}|p{2.4cm}|p{2.4cm}|p{5.6cm}|}
    \hline
    \enc{Decisión} & \enc{Elegida} & \enc{Alternativa} & \enc{Criterio}\\ \hline
    Algoritmo final & Random Forest & Regresión logística, árbol, XGBoost & Mayor PR-AUC en la comparación de la Figura~\ref{fig:comparacion}\\ \hline
    Servicio web & FastAPI & Flask & Validación de tipos y documentación OpenAPI incluidas sin bibliotecas adicionales \cite{ramirez2022deploying}\\ \hline
    Interfaz & React (aplicación independiente) & Streamlit & Frontend desplegable por separado y personalizable; consume la misma API\\ \hline
    Persistencia del modelo & \texttt{joblib} (\textit{Pipeline} completo) & Guardar solo el clasificador & El preprocesamiento viaja con el modelo y se evita el desajuste entre entrenamiento y servicio\\ \hline
  \end{tabular}
\end{table}
